Consider the Linear Algebra (LA) expression $sum((X-UV^T)^2)$ which
defines a typical loss function for approximating a matrix $X$ with a
low-rank matrix $UV^T$. Here, $sum()$ computes the sum of all matrix
entries in its argument, and $A^2$ squares the matrix $A$
element-wise. Suppose $X$ is a sparse, 1M x 500k matrix, and suppose
$U$ and $V$ are dense vectors of dimensions 1M and 500k respectively.
Thus, $UV^T$ is a rank 1 matrix of size 1M x 500k, and computing it
naively requires 0.5 trillion multiplications, plus memory allocation.
Fortunately, the expression is equivalent to
$sum(X^2) - 2U^TXV + U^TU * V^TV$.  Here $U^TXV$ is a scalar that can
be computed efficiently by taking advantage of the sparsity of $X$,
and, similarly, $U^TU$ and $V^TV$ are scalar values requiring only 1M
and 500k multiplications respectively.

Optimization opportunities like this are ubiquitous in machine
learning programs. State-of-the-art optimizing compilers such as
SystemML~\cite{DBLP:reference/bdt/Boehm19},
OptiML\cite{DBLP:conf/icml/SujeethLBRCWAOO11}, and
Cumulon\cite{DBLP:conf/sigmod/HuangB013} commonly implement syntactic
rewrite rules that exploit the algebraic properties of the LA
expressions. For example, SystemML includes a rule that rewrites the
preceding example to a specialized operator \footnote{See the
  \href{https://systemml.apache.org/docs/0.12.0/engine-dev-guide.html}{SystemML
    Engine Developer Guide} for details on the weighted-square loss
  operator \texttt{wsloss}. } to compute the result in a streaming
fashion.  However, such syntactic rules fail on the simplest
variations, for example SystemML fails to optimize $sum((X+UV^T)^2)$,
where we just replaced $-$ with $+$.  Moreover, rules may interact
with each other in complex ways.  In addition, complex ML programs
often have many common subexpressions (CSE), that further interact
with syntactic rules, for example the same expression $UV^T$ may occur
in multiple contexts, each requiring different optimization rules.

%%%% (Dan: I found the next paragraph a bit too vague, so I shorted
%%%% the idea)
% Moreover, these rules resort to heuristics when
% multiple factors impact execution cost. The interaction among these
% factors can become complex: for example, if $UV^T$ in our example
% appears somewhere else in the same program, the original form may be
% preferable since all of its parents can share the computation result,
% and common-subexpression (CSE) elimination avoids recomputing the
% product. The compiler engineer may then add heuristics to only perform
% the rewrite when $UV^T$ does not have other parents. However, if its
% other parents can also be optimized by other rewrites, the heuristics
% will lead the compiler to miss the optimization opportunity. In
% general, heuristics like this complicate the codebase, add maintenance
% cost, and may miss high-impact optimizations.

% These opportunities and challenges call for a more principled approach
% to optimizing linear algebra expressions in machine learning programs.
In this paper we describe \sys, a novel optimization approach for
complex linear algebra programs that leverages relational algebra as
an intermediate representation to completely represent the search
space. SPORES first transforms LA expressions into traditional
Relational Algebra (RA) expressions consisting of joins $*$, union $+$
and aggregates $\sum$.  It then performs a cost-based optimizations on
the resulting Relational Algebra expressions, using only standard
identities in RA.  Finally, the resulting RA expression is converted
back to LA, and executed.

A major advantage of \sys\ is that the optimization rules in RA are
complete.  Linear Algebra seems to require an endless supply of clever
rewrite rules, but, in contrast, by converting to RA, we can prove
that SPORES is complete.  The RA expressions in this paper are
over $K$-relations~\cite{DBLP:conf/pods/GreenKT07}; a tuple $X(i,j)$
is no longer true or false, but has a numerical value, e.g. $5$; in
other words, the RA expressions that result from LA expressions are
interpreted over bags instead of sets.  A folklore theorem states that
two Unions of Conjunctive Queries over bag semantics are equivalent
iff they are isomorphic\footnote{This was claimed, for conjunctive
  queries only, in Theorem 5.2 in~\cite{DBLP:conf/pods/ChaudhuriV93}
  but a proof was never produced; a proof was given for bag-set
  semantics in~\cite{DBLP:journals/tocl/CohenSN05}.  See the
  discussion in~\cite{DBLP:conf/icdt/Green09}.}, which implies that
checking equivalence is decidable.  (In contrast, {\em containment} of
two UCQs with bag semantics is
undecidable~\cite{DBLP:journals/tods/IoannidisR95}; we do not consider
containment in this paper.)  We prove that our optimizer rules are
sufficient to convert any RA expression into its {\em canonical form},
i.e.  to an UCQ, and thus can, in principle, can discover all
equivalent rewritings.

However, we faced a major challenge in trying to exploit the
completeness of the optimizer.  The search space is very large,
typically larger than that encountered in standard database
optimizers, because of the prevalence of unions $+$, large number of
aggregates $\sum$, and frequent common subexpressions. To tackle this,
\sys\ adopts and extends a technique from compilers called {\em
  equality saturation}~\cite{DBLP:journals/corr/abs-1012-1802}.
% to search the space.  Equality saturation 
It uses a data structure called the E-Graph \cite{10.5555/909447} to
compactly represent the space of equivalent expressions, and equality
rules to populate the E-Graph, then leverages constraint solvers to
extract the optimal expression from the E-Graph.  

We have integrated \sys\ into
SystemML~\cite{DBLP:reference/bdt/Boehm19}, and show that it can
derive all hand-coded rules of SystemML.  We evaluated \sys\ on a
spectrum of machine learning tasks, showing competitive performance
improvement compared with more mature heuristic-based optimizers. Our
optimizer rediscovers all optimizations by the latter, and also finds
new optimizations that contribute to speedups of 1.2X to 5X.

We make the following contributions in this paper:
\begin{enumerate}
\item We describe a novel approach for optimizing complex Linear
  Algebra expressions by converting them to Relational Algebra, and
  prove that this approach is complete (Sec.~\ref{sec:la:to:ra}).
\item We present search algorithm based on Equality Saturation that
  can explore a large search space while reusing memory (Sec.~\ref{explore}).
\item We conduct an empirical evaluation of the optimizer using
  several real-world machine learning tasks, and demonstrate it's
  superiority over an heuristics-driven optimizer in SystemML
  (Sec.~\ref{sec:evaluation}).
\end{enumerate}{}

% \dan{Add: ``We make the following contributions in this paper:''
%   Bullet 1, bullet 2, etc}
% \remy{Added above.}
